\section{Contraction of covariance matrices}
\label{sec:contraction}
We now turn to constructing Gaussian matrix product states from local Gaussian
resources~\cite[Secs.~3 and~6]{schuch}. At each site, we distinguish \emph{physical
modes}, which form the final system, from auxiliary \emph{virtual modes}. Neighbouring
sites are joined by EPR measurements on their virtual modes, and we seek the
covariance of the physical modes conditioned on the measurement outcomes.
At the covariance level, this construction leads to a Schur complement.
We first establish the general matrix result, then explain EPR measurements
and apply it to a periodic chain.

\begin{theorem}\label{thm:admissible-positive}
Every admissible $2n\times2n$ covariance matrix is positive definite:
$x^T\gamma x>0$ for every nonzero $x\in\R^{2n}$.
\end{theorem}
\begin{proof}
Put $H=\gamma_{\C}-i\Omega_n$, where the subscript denotes entrywise inclusion
in $\C$. For real $x$, skew-symmetry gives $x^T\Omega_nx=0$, and hence
\[
 x^*Hx=x^T\gamma x\geq0.
\]
For a positive semidefinite Hermitian matrix, a vector has zero quadratic
form precisely when it lies in the kernel. Thus, if $x^T\gamma x=0$, then
$Hx=0$. Taking imaginary parts gives $\Omega_nx=0$. Since $\Omega_n^2=-I$,
we obtain $x=0$. The quadratic form is therefore strictly positive away from
zero. When $n=0$, there is no nonzero vector and the assertion is vacuous.
\end{proof}

\begin{definition}[Contracted covariance]\label{def:contraction}
Let $W$ be a real $(k+2m)\times2n$ matrix. For a real symmetric $2n\times2n$ matrix $\gamma$, write
\[
 V=W\gamma W^T=\begin{pmatrix}E&F\\F^T&G\end{pmatrix},
 \qquad E\in\R^{k\times k}.
\]
The first $k$ output coordinates are to be eliminated; the last $2m$ are
retained in position-then-momentum order. When $E$ is invertible, define
\[
 \gamma_{\mathrm{out}}=G-F^TE^{-1}F.
\]
This is the \emph{Schur complement} of $E$ in $V$. For $k=0$, it reduces to $G$.
\end{definition}
The eliminated coordinates need not form a canonical subsystem. The next
condition makes their symplectic pairings vanish, both among themselves and
with the retained coordinates.

\begin{theorem}\label{thm:contraction}
Let $n,m,k\geq0$ and let $\gamma$ be admissible. Suppose $W^T$ is one-to-one
and
\[
 W\Omega_nW^T=
 \begin{pmatrix}0_{k\times k}&0\\0&\Omega_m\end{pmatrix}.
\]
Then $E$ is positive definite, hence invertible, and
$\gamma_{\mathrm{out}}$ is admissible.
\end{theorem}
\begin{proof}
Theorem~\ref{thm:admissible-positive} and injectivity of $W^T$ give, for $y\ne0$,
\[
 y^TVy=(W^Ty)^T\gamma(W^Ty)>0.
\]
Thus $V$ and its principal block $E$ are positive definite. Congruence of the
input uncertainty matrix yields
\[
 Q=W_{\C}(\gamma_{\C}-i\Omega_n)W_{\C}^*
  =\begin{pmatrix}E_{\C}&F_{\C}\\F_{\C}^*&G_{\C}-i\Omega_m\end{pmatrix}
  \succeq0.
\]
For any complex vector $z$, choose the eliminated component to cancel the first
block row. Direct block multiplication gives
\[
 \begin{pmatrix}-E_{\C}^{-1}F_{\C}z\\z\end{pmatrix}^{\!*}
 Q\begin{pmatrix}-E_{\C}^{-1}F_{\C}z\\z\end{pmatrix}
 =z^*\bigl((G-F^TE^{-1}F)_{\C}-i\Omega_m\bigr)z\geq0.
\]
Here complexification commutes with the inverse of the invertible real matrix
$E$. The output is real symmetric because $E$ and $G$ are symmetric and
$(E^{-1})^T=E^{-1}$. Its uncertainty matrix is consequently Hermitian, so the
quadratic-form inequality proves admissibility. The argument includes empty
eliminated or retained coordinate sets.
\end{proof}

\subsection{EPR measurements and covariance contraction}
Two pairs of canonical quadratures admit the commuting combinations
\[
 X_-=q_A-q_B,\qquad P_+=p_A+p_B,\qquad
 [X_-,P_+]=[q_A,p_A]-[q_B,p_B]=0.
\]
An \emph{EPR measurement} is the joint measurement of these two observables.
Their exact joint eigenvectors are generalized eigenvectors: they do not belong
to the two-mode Hilbert space. In particular, the ideal EPR vector is not a
normalizable state. Finite squeezing gives normalizable approximations, but we
must not treat the infinitely squeezed vector as a density operator of trace one.
Giedke and Cirac~\cite[Secs.~III--IV]{giedke} develop the covariance description of
Gaussian measurements.

Reversing both quadratures of $B$ is a canonical phase rotation. It replaces the
preceding pair by
\begin{equation}\label{eq:epr-plus}
 Y_q=q_A+q_B,\qquad Y_p=p_A-p_B.
\end{equation}
We use this convention for the virtual pairs below. Put
$\Theta=\diag(I_M,-I_M)$ for $M$ such pairs. If we denote the remaining modes
by $C$, the matrix
\[
 L=\begin{pmatrix}I&\Theta&0\\0&0&I_C\end{pmatrix}
\]
maps the quadrature vector in $A,B,C$ order to $(Y,C)$.
The identity $\Theta\Omega_M\Theta^T=-\Omega_M$ gives
$L\Omega L^T=\diag(0,\Omega_C)$, as required in Theorem~\ref{thm:contraction}.

We now compute the corresponding characteristic function, including its
normalization. For a centered Gaussian with covariance $\gamma$, write
\[
 \chi(t)=\langle e^{it^TR}\rangle=\exp(-t^T\gamma t/4).
\]
The factor $1/4$ follows from our convention that $\gamma$ is twice the ordinary
covariance. Thus the joint characteristic function of $Y$ and the retained
quadratures is $\chi(L^T(u,z))$. In $A,B,C$ coordinates,
$L^T(u,z)=(u,\Theta u,z)$: this explains the momentum sign in the contraction
of characteristic-function variables. Write
\[
 L\gamma L^T=\begin{pmatrix}E&F\\F^T&G\end{pmatrix}.
\]
Conditioning on the zero outcome integrates out the Fourier variable $u$.
The quadratic expression in the exponent is
\begin{align*}
 u^TEu+2u^TFz+z^TGz
 &=(u+E^{-1}Fz)^TE(u+E^{-1}Fz)\\
 &\quad+z^T(G-F^TE^{-1}F)z.
\end{align*}
Since $E$ is positive definite, the Gaussian integral is finite. Dividing by
its value at $z=0$ leaves the characteristic function with covariance
$G-F^TE^{-1}F$. More generally, for means $d_Y,d_C$ and outcome $a$, the usual
Gaussian conditioning formula is
\begin{equation}\label{eq:conditional-mean}
 d_{C\mid a}=d_C+F^TE^{-1}(a-d_Y),\qquad
 \gamma_{C\mid a}=G-F^TE^{-1}F.
\end{equation}
The outcome changes the mean, but not the conditional covariance. We understand
conditioning here through the outcome density, not as an event of positive
probability at one exact value. These formulas give the measurement interpretation;
Theorem~\ref{thm:contraction} establishes the covariance inequality directly.

In the Choi--Jamio{\l}kowski construction, the resource associated with a map
is obtained formally by applying the map to one half of an ideal EPR pair.
EPR contraction with the other half represents feeding an input to the map.
This turns the \revs{construction of Gaussian matrix product states} into contractions of local resources, but
still involves ideal EPR measurements. Moreover, a channel representative obtained
from an infinitely squeezed pair need not be a finite covariance matrix: the
identity channel already illustrates the difficulty. The results below use
finite admissible local matrices. Schuch, Wolf and Cirac describe the resource
construction in Sections~3 and~6 of~\cite{schuch};
their Section 5 replaces infinitely entangled bonds by finite ones while moving
the infinite entanglement into the local maps.

\subsection{A periodic chain}
At each site $i\in\mathbb Z/N\mathbb Z$, we introduce $M$ left virtual modes $A_i$,
$M$ right virtual modes $B_i$, and one retained mode $C_i$.
The pair on each edge joins $B_{i-1}$ to $A_i$. We assume $N\geq3$ and $M\geq1$.
The integer $M$ counts the \emph{virtual modes per bond}. Each bosonic mode has
an infinite-dimensional Hilbert space, so even $M=1$ can support entanglement.
With no virtual modes, $M=0$, the virtual Hilbert space is one-dimensional;
this is the counterpart of bond dimension $D=1$ for a qudit MPS. We can also
represent product covariances with $M\geq1$ by adding unused virtual modes.
A bound on $M$ alone does not bound entanglement entropy: the squeezing must
also be controlled.

\begin{definition}[Periodic covariance]\label{def:periodic}
Put $d=2M+1$ and let $\Gamma_i$ be real $2d\times2d$ matrices. The local mode
order is $(A_{i,1},\ldots,A_{i,M},B_{i,1},\ldots,B_{i,M},C_i)$, with positions
preceding momenta. Define $\gamma_{\rm loc}$ on $Nd$ modes by
\[
 (\gamma_{\rm loc})_{(u,i,a),(v,j,b)}=
 \begin{cases}(\Gamma_i)_{(u,a),(v,b)},&i=j,\\0,&i\ne j,\end{cases}
 \qquad u,v\in\{q,p\}.
\]
Globally, positions precede momenta; within either group, sites precede local
modes. The real linear map $W:\R^{2Nd}\to\R^{2NM}\oplus\R^{2N}$ has eliminated
coordinates
\[
 (Wx)_{q,i,a}=x_{q,A_i,a}+x_{q,B_{i-1},a},\qquad
 (Wx)_{p,i,a}=x_{p,A_i,a}-x_{p,B_{i-1},a},
\]
and retained coordinates $(Wx)_{q,C_i}=x_{q,C_i}$ and
$(Wx)_{p,C_i}=x_{p,C_i}$. Each output group is in position-then-momentum order.
When its eliminated block is invertible, the \emph{periodic covariance} is
\[
 \gamma_{\rm per}=\mathrm{SC}_{\R^{2NM}}(W\gamma_{\rm loc}W^T),
\]
where $\mathrm{SC}$ denotes the Schur complement of Definition~\ref{def:contraction}.
\end{definition}
For comparison with~\revs{\cite{schuch}}, we group all $A$ coordinates together, then all $B$
and all $C$ coordinates. If $R$ shifts the site index to the right,
$(Rx)_i=x_{i-1}$, the same map is
\[
 W=\begin{pmatrix}I_A&R\Theta_B&0\\0&0&I_C\end{pmatrix}.
\]
The shift acts identically on positions and momenta and preserves each mode
label. This is the matrix in Section~6, Eq.~(3) of~\cite{schuch}.

\begin{theorem}\label{thm:periodic}
Suppose every $\Gamma_i$ is admissible. Then $W^T$ is injective,
\[
 W\Omega_{Nd}W^T=\diag(0_{2NM},\Omega_N),
\]
the eliminated block of $W\gamma_{\rm loc}W^T$ is positive definite, and
$\gamma_{\rm per}$ is admissible.
\end{theorem}
\begin{proof}
The independent sum of the local covariances is admissible by
Theorem~\ref{thm:sum}, iterated over the sites. Reordering its positions and
momenta together gives $\gamma_{\rm loc}$ and preserves the uncertainty inequality.
For an output vector $(y_q,y_p,z_q,z_p)$, the $A$ components of its image under
$W^T$ are $y_q,y_p$, and the $C$ components are $z_q,z_p$. Hence $W^T$ is injective.
Each eliminated position--momentum pairing has an $A$ contribution $+1$ and a
$B$ contribution $-1$. They cancel, including after the common cyclic shift.
Distinct mode pairs and all pairings with $C$ vanish. The retained pairings
form $\Omega_N$. Apply Theorem~\ref{thm:contraction}.
\end{proof}

\subsection{Three correlated sites}
\begin{proposition}\label{prop:three-site}
Take $N=3$, $M=1$, and at each site let
\[
 K=\begin{pmatrix}0&0&1\\0&0&1\\1&1&0\end{pmatrix},\qquad
 S=\begin{pmatrix}I_3&0\\K&I_3\end{pmatrix},\qquad \Gamma_i=SS^T.
\]
Then $S$ is symplectic, $\Gamma_i$ is admissible, and the periodic covariance is
\[
 \gamma_{\rm per}=\diag(Q,P),\qquad
 Q=\begin{pmatrix}3/5&1/5&1/5\\1/5&3/5&1/5\\1/5&1/5&3/5\end{pmatrix},\quad
 P=\begin{pmatrix}2&-1/2&-1/2\\-1/2&2&-1/2\\-1/2&-1/2&2\end{pmatrix}.
\]
It is admissible. In particular, distinct sites have position--position covariance
entry $1/5$ and momentum--momentum covariance entry $-1/2$.
\end{proposition}
\begin{proof}
Since $K=K^T$, direct multiplication gives $S\Omega_3S^T=\Omega_3$.
The covariance $I_6$ is an independent sum of three one-mode vacua, so
Theorem~\ref{thm:congruence} gives admissibility of $SS^T$.
Let $R$ be the three-site cyclic right shift and $U$ the $3\times3$ all-ones
matrix. Multiplying the local covariance blocks gives
\[
 E=\begin{pmatrix}2I&0\\0&4I-R-R^T\end{pmatrix},\quad
 F=\begin{pmatrix}0&I+R\\I-R&0\end{pmatrix},\quad
 G=\begin{pmatrix}I&0\\0&3I\end{pmatrix}.
\]
Here $R+R^T=U-I$ and $U^2=3U$. Consequently
\[
 (4I-R-R^T)^{-1}=(5I-U)^{-1}=\tfrac15 I+\tfrac1{10}U,
\]
as multiplication checks. Since $U(I-R)=0$, the two diagonal blocks of
the Schur complement are
\begin{align*}
 Q&=I-(I-R)^T(\tfrac15I+\tfrac1{10}U)(I-R)
   =\tfrac25I+\tfrac15U,\\
 P&=3I-\tfrac12(I+R)^T(I+R)=\tfrac52I-\tfrac12U.
\end{align*}
The off-diagonal blocks vanish. Theorem~\ref{thm:periodic} proves admissibility.
\end{proof}

For this example, $QP=I_3$, so the associated Gaussian state is pure. Its
single-site covariance is $\diag(3/5,2)$, with determinant $6/5>1$; the
single-site reduced state is therefore mixed. Thus the state is entangled
across every partition of one site from the other two. These deductions use
the Gaussian purity criterion~\cite[Secs.~II.C--D]{weedbrook}\revs{, which is cited, not
proved here; Proposition~\ref{prop:three-site} establishes the displayed covariance and its admissibility}.